\section{古诗文赏析的跨任务微调}

为观察同一预训练主干对另一任务的适应，本研究另行微调诗文赏析模型。该分支直接继承原始GPT，以解释性段落为目标，使用独立权重。

语料来自\href{https://huggingface.co/datasets/byj233/ChinesePoetryLibrary}{ChinesePoetryLibrary}和\href{https://huggingface.co/datasets/Papersnake/gushiwen}{Papersnake/gushiwen}，保留发布者的AI赏析来源标记。按作品与赏析分组，清理后训练798,397条，验证和测试各8,072条；超长、异常字形与精确重复记录合计过滤2,647条。训练使用完整正文，词表扩至15,829、上下文扩至1,536，原字符和前128个位置继承，总参数9,179,904。

输入包含诗文、任务说明和程序字音依据，仅监督赏析正文及EOS。有效批量128，AdamW峰值学习率$1.5\times10^{-4}$，预热3\%后余弦下降，训练一轮共6,238步、221,304,997个目标词元，约32.2分钟，按验证损失选择第6,236步权重。字音依据来自\href{https://github.com/charlesix59/chinese_word_rhyme}{平水韵字表}，展示时与模型正文分开。

表~\ref{tab:appreciation}采用相同输入比较任务初始化与微调后模型：交叉熵使用1,024条测试，自由生成使用固定64条，均为贪心解码。初始化组采用相同扩展词表与位置表，指标排除程序字音首句。

\begin{table}[H]
\centering\small
\caption{诗文赏析任务微调前后对比}\label{tab:appreciation}
\begin{tabularx}{\linewidth}{Yrr}
\toprule 指标&预训练初始化&监督微调后\\\midrule
赏析正文交叉熵&7.5196&1.5715\\
自然结束率&73.4\%&98.4\%\\
严重四字符循环率&67.2\%&20.3\%\\
生成chrF&1.7933&20.7481\\\bottomrule
\end{tabularx}
\end{table}

以明代李梦阳《十五夜》为例，原文从“去年燕山南、易水北”写到“今年有月复有台”，两组实际输出如下。

\modeltext{微调前（节选）}{我心如此意，不可怜我心如此意不可怜我心如此意，不见。}

\modeltext{微调后（节选）}{从文学手法上看，这首诗运用了对比的手法，如‘去年燕山南、易水北，中宵待月不可得’，将过去的时光与如今的‘今年有月复有台，海门烟灭没’的对比，突出时光的流逝和世事的变迁。}

微调后开始围绕原文中的时间对比组织解释。另一方面，樊增祥《踏莎行》的输出将“女儿心事莲花净”与同一句话解释为“对比”。这些实例呈现了任务体裁的变化与修辞判断错误，完整对照见附录B。
